\section{Experiments}
\label{sec:experiments}

We evaluate RQ2: how well next-phase elimination hazard can be predicted from
the team-phase features alone, and whether the dynamic team graph adds signal
beyond them. All experiments in this section use the full graph-eligible
population described in Table~\ref{tab:data-funnel}. The train/validation/test
split is performed at the match level with seed 42, producing 17,460 training
matches, 2,182 validation matches, and 2,182 test matches. The split is
random rather than chronological and is not stratified; mode proportions
nevertheless balance across splits (squad 69.7\%, 68.7\%, and 70.9\% of
train, validation, and test). Because team-phase cells within a match are
correlated, every paired test resamples matches, never rows.

\paragraph{Models.}
We compare five models. The tabular baseline is a histogram gradient
boosting classifier trained on a 72-column tabular view of the same team-phase
rows; the view additionally includes raw phase timing and categorical map and
mode identifiers, so the baseline sees match context that the graph nodes do
not. It uses the same leakage-safe target definition as the graph models but
ignores adjacency tensors. The no-graph DynamicSTGNN uses the same neural architecture
with all edge channels disabled; it is therefore a sequence-only neural
baseline over team-phase node features. The proximity model enables only
spatial-nearness edges. The proximity+combat model adds directed combat edges.
The full DynamicSTGNN uses all three edge channels: proximity, combat, and
co-pressure.

\paragraph{Representation Conditions.}
To separate the representation from the model class, we refit the same
gradient-boosting pipeline on nested feature sets: prior-work-style strategy
scores (five features: phase index, normalized aggression and activity scores,
a derived strategy type, and that type's historical survival rate, following
the phase-based scoring of \cite{sagong2025hunterCamper}); the twelve raw
quantities behind those scores; the STDB movement and exposure features alone;
conventional telemetry aggregates; and the conventional set plus the STDB
features---the 72-column input of the tabular baseline. Refitting that last
set reproduces the Table~\ref{tab:main-hazard-results} baseline to within
0.0001 AUC. All conditions use the same split and aligned test cells.

\paragraph{Training Protocol.}
DynamicSTGNN models are trained for 12 epochs with batch size 8, AdamW
optimization (learning rate \(10^{-3}\), weight decay \(10^{-4}\)) with cosine
annealing and gradient-norm clipping, and validation loss for checkpoint
selection. The loss is masked binary cross-entropy for the hazard target.
The positive rate is 0.32 and training uses no reweighting (hazard positive
weight 1). Hyperparameters were fixed from preliminary runs; no per-model
search was performed. Test metrics are computed from the best validation
checkpoint.

\paragraph{Aligned Evaluation.}
Because tabular and graph pipelines can differ slightly in which team-phase
cells are materialized, we report the main comparison on the common aligned
test cells shared by all compared models. This produces 134,245 target-defined
team-phase cells from 2,182 test matches. Statistical tests use one match as
one paired unit. For each model and match we compute BCE, AUC, and AP, then use
one-sided Wilcoxon signed-rank tests where positive improvement means the left
model is better.

\paragraph{Population Consistency.}
All tables in Sections~\ref{sec:experiments} and~\ref{sec:results} use the
same full-DB split and aligned test cells. Earlier exploratory seed-repeat and
random-edge sanity experiments were run on older snapshots of the dataset; we
do not mix those values into the main result tables unless explicitly labeled
as auxiliary evidence.

\paragraph{Metrics.}
We report ROC-AUC and average precision as primary metrics. ROC-AUC summarizes
how often a target-positive team-phase cell receives a higher predicted hazard
than a target-negative cell; higher is better. Average precision is the area
under the precision-recall curve and is useful because hazard positives are
not uniformly balanced across phases; higher is better. We also report
probability diagnostics: binary cross-entropy (BCE), Brier score, and expected
calibration error (ECE). BCE is the negative log-likelihood of the binary
hazard label under the predicted probability, and Brier score is the mean
squared probability error. ECE partitions predictions into 10 equal-width
probability bins and averages the absolute gap between each bin's mean
predicted hazard and observed elimination rate, weighted by bin size. BCE,
Brier, and ECE are lower-is-better. Probability quality matters alongside
ranking quality because the task is a hazard prediction problem. We
additionally report mode-wise discrimination to check that gains are not a
mode shortcut.
